
\medterm{Waardenburg Syndrome} A group of inherited conditions that may cause hearing loss and unusual coloring of the eyes, hair, or skin. Some forms also affect the arms or the nerves controlling the bowel; findings vary by type.


\medterm{WAGR Syndrome} A rare genetic condition that may cause an increased risk of Wilms kidney tumor, missing or poorly formed irises, urinary or reproductive differences, and developmental delay. Children need regular specialist monitoring, especially for kidney tumors.

\medterm{Waist Circumference} The distance around the abdomen, measured at a standard level. A larger measurement can indicate more abdominal fat and a higher risk of diabetes and heart disease, but it must be interpreted with other health information.

\medterm{Waist-Hip Ratio} A measurement comparing waist size with hip size. It provides an estimate of where body fat is stored and may add information about diabetes and cardiovascular risk, but it does not diagnose disease by itself.

\medterm{Wait List} An official list of people waiting for an organ transplant. Placement and priority depend on medical urgency, expected benefit, compatibility, waiting time, and the rules of the transplant system.

\medterm{Wakefulness} The state of being conscious and able to respond to the surroundings. Sleep, medicines, illness, brain injury, and metabolic problems can change the level of wakefulness.

\medterm{Waldenstrom Macroglobulinemia} A rare blood cancer in which abnormal immune cells build up in the bone marrow and make too much of an antibody called IgM. It may cause anemia, bleeding, vision or nerve problems, enlarged lymph nodes, or no symptoms at first.

\synonyms
Lymphoplasmacytic Lymphoma

\medterm{Waldenstrom'S Macroglobulinemia} A lymphoplasmacytic cancer of the bone marrow that produces excess IgM and may cause
anemia, enlarged lymph nodes, or hyperviscosity symptoms.

\medterm{WaldenströM'S Macroglobulinemia} A slow-growing blood cancer in which abnormal immune cells make too much of an antibody called IgM. The blood can become unusually thick, causing blurred vision, headaches, bleeding, confusion, or numbness.

\medterm{Walker-Warburg Syndrome} A rare inherited disorder causing severe muscle weakness together with major abnormalities of the brain and eyes from birth. Babies may have low muscle tone, seizures, poor feeding, enlarged fluid spaces in the brain, and abnormal eye development; the condition is usually life-threatening in infancy.

\synonyms
WWS; Muscle-Eye-Brain Disease (Severe Type); HARD +/- E Syndrome

\medterm{Walking Abnormalities} Walking abnormalities include changes in speed, balance, stride, posture, or coordination caused by pain, weakness, sensory loss, joint disease, or neurologic illness.

\medterm{Walking Speed} The rate at which a person walks, often measured over a defined distance. It can provide information about mobility and function, but no single speed determines disability in every person.

\medterm{Wall Stress} The strain placed on the wall of a heart chamber by the pressure inside it and the chamber’s size. High wall stress makes pumping harder and may contribute to heart enlargement or failure.

\medterm{Wallerian Degeneration} Breakdown of the part of a nerve beyond a severe injury after it loses contact with the nerve cell. The nerve may regrow slowly if the surrounding pathway remains intact, but recovery depends on the injury.


\medterm{Walleyed} Walleyed describes eyes that appear misaligned or divergent; clinically, the condition may be strabismus and requires eye assessment when persistent.

\medterm{Wandering Spleen} A spleen that can move unusually because its supporting bands are loose or missing. It may cause no symptoms or lead to abdominal pain, enlargement, twisting, loss of blood supply, or sudden severe illness.


\medterm{Warfarin} A blood-thinning medicine that reduces production of vitamin-K-dependent clotting proteins. It prevents and treats clots but requires regular blood tests because food, illnesses, and many medicines alter its effect.

\medterm{Warfarin Syndrome} A pattern of fetal abnormalities caused by warfarin exposure during pregnancy, including nasal and skeletal abnormalities, eye or neurologic defects, and fetal or neonatal bleeding.

\medterm{Warm Autoimmune Hemolytic Anemia} An illness in which the immune system makes antibodies that cause red blood cells to be destroyed too early, mainly in the spleen. It can cause tiredness, pale or yellow skin, dark urine, shortness of breath, or a fast heartbeat and requires medical evaluation.

\synonyms
wAIHA; Warm antibody autoimmune hemolytic anemia

\medterm{Warm Ischemia} Lack of blood and oxygen reaching tissue while it remains at body temperature. It can quickly cause cell damage and is important during some injuries, operations, and organ transplantation.

\medterm{Wart} A usually harmless rough skin growth caused by human papillomavirus. Warts may spread by contact and often disappear without treatment, but painful, bleeding, changing, or uncertain growths should be examined.

\synonyms
Verruca

\medterm{Wart Remover Poisoning} Toxic exposure to a wart-removal product, often involving salicylic acid or another keratolytic. Severity depends on the substance, amount, route, symptoms, and duration of exposure.

\medterm{Warthin Tumor} A usually harmless tumor of a salivary gland, most often the parotid near the ear. It is strongly associated with tobacco use and may occur on both sides; evaluation distinguishes it from other growths.

\medterm{Warty Dyskeratoma} A usually harmless skin growth, often on the head or neck, with a small central dip filled with keratin. It may look like skin cancer, so examination and sometimes a biopsy are needed.



\medterm{Wasting} Abnormal loss of body weight, muscle, or other tissue, often associated with chronic disease, inadequate nutrition, inflammation, or cancer. Severe wasting may impair strength, immunity, and recovery.

\medterm{Watchful Waiting} Planned observation without immediate active treatment, with clinical review and treatment
if symptoms develop or the condition progresses. It may be used for selected cancers and other slowly progressive
conditions.

\medterm{Water Birth (Underwater Delivery)} Delivery of a baby while the birthing person remains in a tub or pool of warm water.
It may reduce perceived labor pain for some people, but evidence about delivery-related
benefits and risks is limited. Eligibility, monitoring, and facility policies determine whether
it is offered.

\synonyms
Underwater Delivery

\medterm{Water Blister} A fluid-filled blister, usually containing clear serum, caused by friction, pressure, burns, or some skin disorders. The blister should be protected from further trauma and assessed if infected or widespread.

\medterm{Water Brash} A sudden rise of watery or salty-tasting saliva into the mouth, usually because stomach acid has refluxed into the esophagus. It may occur with heartburn, sour taste, throat irritation, or regurgitation.

\medterm{Water Channel} A protein pore in a cell membrane that lets water cross quickly. These pores help the kidneys and other tissues move water and maintain the body’s fluid balance.

\medterm{Water Consumption} Intake of water and other fluids to replace normal losses and maintain hydration. Requirements vary with body size, diet, activity, climate, pregnancy, illness, and kidney or heart function; thirst and urine concentration can help guide intake, while excessive intake can cause dangerous hyponatremia.


\medterm{Water Intoxication} A dangerous fall in blood sodium caused by drinking or retaining more water than the body can remove. Severe cases can cause headache, confusion, seizures, or coma and need urgent treatment.

\medterm{Water On The Brain} A nontechnical term usually referring to hydrocephalus, an abnormal accumulation or impaired circulation of cerebrospinal fluid that enlarges the ventricles and may increase intracranial pressure.

\medterm{Water On The Knee} Alternate index label; see \textit{Swollen knee}.


\medterm{Water Retention} Excess accumulation of fluid in the body's tissues or cavities, commonly causing swelling and sometimes resulting from heart, kidney, liver, venous, hormonal, or medication-related disorders.



\medterm{Water-Borne Disease} An illness acquired through water contaminated with infectious organisms or harmful chemicals. Examples include cholera, typhoid fever, and hepatitis A; prevention includes safe drinking water, sanitation, hand hygiene, and appropriate treatment of contaminated water.

\medterm{Water-Deprivation Test} A supervised test that briefly limits drinking while urine concentration and blood results are monitored. It helps identify why a person produces excessive amounts of dilute urine and must be done under medical supervision because dehydration and high sodium can occur.

\medterm{Water-Hammer Pulse (Corrigan Pulse)} A pulse that feels forceful at first and then suddenly collapses. It is often associated with severe leakage of the aortic valve or other causes of a wide pulse pressure.

\synonyms
Corrigan Pulse

\medterm{Water-Soluble Vitamin} A vitamin that dissolves in water and is stored only in limited amounts. The group includes vitamin C and the B vitamins, so regular intake is generally needed.


\medterm{Waterborne Bacterial Disease} An infection acquired from water contaminated with bacteria, such as cholera, typhoid fever, or leptospirosis; prevention depends on safe water, sanitation, and food hygiene.

\medterm{Waterhouse-Friderichsen Syndrome} Waterhouse-Friderichsen syndrome is adrenal hemorrhage and acute adrenal failure associated with overwhelming sepsis, classically meningococcal infection, causing shock and coagulopathy.

\medterm{Watery Eyes} Excessive tearing caused by irritation, allergy, infection, dry-eye reflex tearing, or
blockage of the tear-drainage system. Persistent unilateral tearing or pain requires
ocular assessment.

\medterm{Wax} A water-resistant substance produced by ceruminous glands in the ear canal; earwax protects the canal and usually clears naturally.

\medterm{Wax Poisoning} Toxic exposure to a wax product, usually causing little harm after a small accidental ingestion but potentially causing choking, aspiration, irritation, or toxicity when mixed with other chemicals.

\medterm{Waxy Casts} Large, smooth, pale tube-shaped particles seen under a microscope in urine. They form in damaged kidney tubules and can suggest severe or long-standing kidney disease, especially when other urine or blood tests are abnormal.

\medterm{WBC Count} A measurement of the number of white blood cells in a volume of blood, used with the differential count to evaluate infection, inflammation, immune disorders, marrow disease, or treatment effects.

\medterm{Weakness} Reduced ability to generate force, which may be generalized or limited to one body region. Causes include disorders of muscle, nerve, neuromuscular junction, brain, or spinal cord, as well as metabolic illness, infection, medication effects, and deconditioning; sudden one-sided weakness or weakness with breathing, swallowing, or speech difficulty requires emergency assessment.

\medterm{Weakness, Generalized} Reduced strength affecting multiple body regions, caused by systemic illness, metabolic or electrolyte disturbance, medication, infection, muscle disease, nerve disease, neuromuscular-junction disease, or brain or spinal-cord disease. Sudden or progressive weakness, especially with speech, facial, breathing, or walking difficulty, requires urgent assessment.

\medterm{Weaning} Gradual reduction and eventual discontinuation of breastfeeding, bottle-feeding, or a medicine or other treatment, according to the relevant clinical context.

\medterm{Wear Debris} Tiny particles produced when an artificial joint or other implant wears down. The particles can trigger inflammation, damage nearby bone, loosen the implant, and eventually cause it to fail.

\medterm{Weaver Syndrome} A rare genetic condition causing rapid growth before and after birth, a large head, distinctive facial features, low muscle tone, and delayed development or learning. Bone age may be advanced. Most cases involve a change in the \textit{EZH2} gene; care is tailored to the child’s needs.

\medterm{Weaver'S Bottom} Inflammation of the ischial bursa over the sitting bone, causing pain in the buttock that may worsen with prolonged sitting. It is also called ischial bursitis.


\medterm{Webbed Fingers} Partial or complete joining of adjacent fingers by skin or soft tissue, known as syndactyly; it may be congenital or acquired.

\medterm{Webbed Neck} Extra folds of skin extending from the sides of the neck toward the shoulders. They may be present from birth and can occur with conditions such as Turner syndrome or Noonan syndrome.

\medterm{Webbing Of The Fingers Or Toes} Partial or complete joining of adjacent fingers or toes by skin or other tissue. It is called syndactyly and may be congenital or acquired.

\medterm{Weber} A name used in hearing assessment, especially the Weber test, in which a vibrating tuning fork is placed on the forehead or skull. It helps compare hearing between the two ears but is not a unit of measurement.

\medterm{Weber Test} A quick hearing test using a vibrating tuning fork placed on the middle of the forehead or skull. The examiner asks which ear hears the sound louder; the result helps distinguish blockage-related hearing loss from inner-ear or nerve-related loss.

\medterm{Wedge Excision} Surgical removal of a triangular or wedge-shaped portion of tissue; commonly performed for breast biopsies, cervical lesions, and skin tumors.

\medterm{Wedge Resection} Surgery that removes a small wedge-shaped piece of tissue, commonly from the lung or breast, to diagnose or treat a localized lesion. It removes less tissue than removal of an entire lobe.


\medterm{Weever Fish} A venomous marine fish whose dorsal spines can cause a painful sting, local swelling, and occasionally systemic symptoms.

\medterm{Wegener'S Granulomatosis} Alternate historical term; see \textit{Granulomatosis With Polyangiitis}.

\medterm{Wegeners Granulomatosis (Granulomatosis With Polyangiitis)} Alternate index label for Granulomatosis With Polyangiitis.

\medterm{Weight Gain} Increase in body weight; may be intentional (muscle building, pregnancy) or unintentional (fluid retention, endocrine disorders, medication effects).

\medterm{Weight Gain In Pregnancy} Recommended weight gain depends on weight before pregnancy and whether one or more babies are expected. For one baby, usual ranges are about 12.5--18 kg if underweight, 11.5--16 kg at usual weight, 7--11.5 kg if overweight, and 5--9 kg with obesity; prenatal clinicians monitor the pattern and fetal growth.

\medterm{Weight Gain – Unintentional} Increase in body weight without a planned change in diet or activity, often caused by fluid retention, medicines, endocrine disease, reduced activity, or changes in eating; rapid swelling or breathlessness needs prompt assessment.

\medterm{Weight Loss} A decrease in body weight that may be intentional or unintentional. Intentional loss is usually achieved through sustainable changes in eating, physical activity, and behavior; unintentional or rapid loss can indicate illness, inadequate nutrition, or medication effects and may require medical evaluation.

\medterm{Weight Loss - Unintentional} Unintentional weight loss is reduction in body weight without planned dietary or activity change and may result from malignancy, endocrine disease, infection, gastrointestinal disease, medication, or mood disorder.

\medterm{Weight Loss Medications} Medicines prescribed for selected people with obesity to help reduce weight, usually together with nutrition and activity changes. Choice depends on health conditions, other medicines, side effects, cost, and treatment goals.
\medterm{Weight-Loss Medicines} Alternate index label; see \textit{Weight Loss Medications}.

\medterm{Weight-Reducing Diets} Dietary plans intended to create a sustained energy deficit and reduce excess body weight, ideally individualized with attention to nutrition, health conditions, and long-term adherence.


\medterm{Weil'S Disease} A zoonotic bacterial infection caused by spirochaetes of the genus Leptospira,
transmitted through contact with water or soil contaminated with the urine of infected
animals (rodents, cattle, dogs). See Leptospirosis.

\medterm{Weils Disease (Leptospirosis)} Alternate index label for Leptospirosis.

\medterm{Wellens Syndrome} A specific pattern on a heart tracing that can signal severe narrowing of a major heart artery, even after chest pain stops. It indicates high heart-attack risk and needs urgent specialist assessment; exercise testing can be dangerous.

\synonyms
Wellens Pattern; LAD Coronary T-Wave Syndrome; Wellens Sign

\medterm{Wells Score For Pulmonary Embolism} A clinical score estimating the likelihood of a blood clot in the lungs. It considers leg-clot signs, heart rate, recent surgery or immobility, previous clots, coughing blood, cancer, and alternative diagnoses; the result helps decide whether blood testing or lung imaging is needed.

\synonyms
Wells Criteria for PE; Wells Score (PE)

\medterm{Welt} A raised, often itchy area of skin caused by localized swelling, commonly seen in hives or after an insect bite. A welt may fade within hours, but persistent or severe reactions need evaluation.

\medterm{Wen (Epidermoid Cyst)} A common, usually harmless lump beneath the skin filled with a soft protein called keratin. It may become swollen, painful, or infected and can be removed if necessary.

\synonyms
Epidermoid Cyst

\medterm{Wenckebach Phenomenon (Mobitz Type I AV Block)} A heart-rhythm pattern in which each heartbeat takes slightly longer to pass from the upper to lower chambers until one beat is missed. It may be harmless or reflect heart disease, depending on the setting.

\synonyms
Mobitz Type I AV Block

\medterm{Wens} Alternate index label; see \textit{Wen (Epidermoid Cyst)}.

\medterm{Werdnig-Hoffmann Disease} The most severe early form of inherited spinal muscular atrophy, also called type 1 spinal muscular atrophy. It usually begins before 6 months of age and causes very weak muscles, poor head control, reduced reflexes, a weak cry, feeding difficulty, and breathing problems. Babies generally cannot sit without support. New treatments can improve survival and function, but affected infants need urgent specialist care and breathing and feeding support.

\synonyms
Spinal Muscular Atrophy Type 1 (SMA 1); Severe Infantile Spinal Muscular Atrophy

\medterm{Wermer Syndrome} Alternate index label; see \textit{Multiple endocrine neoplasia, type 1 (MEN 1)}.

\medterm{Wermer'S Syndrome} Alternate index label; see \textit{Multiple endocrine neoplasia, type 1 (MEN 1)}.

\medterm{Werner Syndrome} A rare inherited condition that causes signs of aging earlier than usual, often after puberty. Features may include short stature, early cataracts, thin or tight skin, gray hair, diabetes, weak bones, and increased cancer risk. Regular medical monitoring and treatment of each complication are needed.

\medterm{Werner'S Syndrome} A rare inherited condition that causes signs of aging earlier than usual, often after puberty. It may cause early cataracts, short stature, thin skin, gray hair, diabetes, weak bones, and increased cancer risk.

\medterm{Wernicke Area} A brain region, usually in the left temporal lobe, that helps a person understand spoken and written language. Damage can cause fluent speech that lacks meaning and poor understanding of language.

\medterm{Wernicke Encephalopathy} A sudden brain disorder caused by severe vitamin B1 deficiency. It may cause confusion, difficulty walking, and abnormal eye movements and can lead to permanent damage without prompt treatment.

\medterm{Wernicke'S Encephalopathy} An acute neurologic disorder caused by thiamine deficiency, classically involving confusion,
eye-movement abnormalities, and gait or coordination impairment. It requires prompt
thiamine treatment.

\medterm{Wernicke-Korsakoff Syndrome} A brain disorder caused by severe vitamin B1 deficiency. It may begin with confusion and poor coordination and leave lasting memory problems, difficulty learning, and inaccurate made-up explanations; alcohol misuse, poor nutrition, and prolonged vomiting are common causes.

\medterm{Wernicke’S Area} A brain region important for understanding spoken and written language.

\synonyms
Wernicke area


\medterm{West Nile Encephalitis} Inflammation of the brain caused by neuroinvasive West Nile virus, transmitted mainly by mosquitoes. It may cause
fever, headache, confusion, tremor, weakness, seizures, or coma. There is no specific routine antiviral treatment;
severe illness requires hospital-based supportive care.

\medterm{West Nile Virus} A mosquito-borne flavivirus that can cause asymptomatic infection, a febrile illness, or neuroinvasive
disease such as meningitis, encephalitis, or acute flaccid paralysis. Risk of severe disease
is higher in older adults and people with certain medical conditions.

\medterm{West Nile Virus Infection} West Nile virus infection is a mosquito-borne illness ranging from fever and body aches to neuroinvasive meningitis, encephalitis, weakness, or paralysis.

\medterm{West Syndrome} A serious seizure disorder beginning in infancy, with repeated sudden spasms and an abnormal brain-wave pattern. Development may slow or regress; causes include brain injury, genetic conditions, and metabolic disease, and urgent treatment is needed.

\medterm{Westermark Sign} An area of unusually reduced blood-vessel markings on a chest X-ray, sometimes seen beyond a clot in a lung artery. It may suggest pulmonary embolism but is not reliable enough to diagnose it alone.


\medterm{Westley Croup Score} A clinical score used to estimate how severely croup is obstructing a child's airway. It considers noisy breathing, chest retractions, air movement, bluish coloring, and alertness; a higher score suggests more severe illness and helps guide urgent treatment and observation.

\synonyms
Westley Clinical Scoring System for Croup; Croup Severity Score

\medterm{Wet Age-Related Macular Degeneration} Alternate index label; see \textit{Wet macular degeneration}.

\medterm{Wet AMD} Alternate index label; see \textit{Wet macular degeneration}.

\medterm{Wet Dream} Involuntary ejaculation during sleep, usually occurring in adolescence or adulthood and generally a normal physiological event.

\medterm{Wet Dry Drowning (Drowning)} Alternate index label for Drowning.

\medterm{Wet Gangrene} Tissue death complicated by bacterial infection and tissue swelling, producing moist, discolored, rapidly spreading necrosis. It requires urgent medical assessment.

\medterm{Wet Macular Degeneration} The neovascular form of age-related macular degeneration in which abnormal blood vessels grow beneath or within the retina and leak fluid or blood, causing rapid central-vision loss.

\synonyms
Macular Degeneration, Wet
Wet Age-Related Macular Degeneration
Wet AMD

\medterm{Wharton Jelly} The soft, jelly-like tissue surrounding the blood vessels in the umbilical cord. It cushions the vessels and helps prevent them from being squeezed or bent during pregnancy and labor.

\medterm{Wharton'S Jelly} Gelatinous connective tissue within the umbilical cord that surrounds and supports the umbilical blood vessels and helps protect them from compression.

\medterm{Wheal} A temporary, raised, often itchy area of skin caused by swelling in the skin. It commonly occurs with hives, an allergic reaction, or an insect bite and may move or disappear within hours.

\medterm{Wheal And Flare} A temporary raised, swollen area of skin surrounded by redness. It is usually caused by release of histamine during hives, an allergy, or another immediate skin reaction.

\medterm{Wheat Allergy} An immune reaction to proteins in wheat that may cause hives, stomach symptoms, vomiting, or breathing problems.

\synonyms
Allergy, Wheat
\medterm{Wheeled Stretcher} Mobile patient transport device with wheels; used in hospitals and emergency services for moving patients who cannot walk.

\medterm{Wheeze} A high-pitched whistling sound caused by narrowed airways. It often occurs with asthma, infection, allergy, or chronic lung disease; sudden wheezing after choking or with severe breathlessness needs urgent care.

\medterm{Wheezing} A musical or whistling sound caused by narrowed lower airways, usually louder when breathing out. It may result from asthma, infection, chronic lung disease, allergy, airway swelling, mucus, or an inhaled object.

\medterm{Whiplash} A neck injury caused by the head moving suddenly backward and forward, often in a vehicle crash but also after a fall or sports injury. Pain, stiffness, headache, dizziness, or arm tingling may appear immediately or later. Most cases improve over weeks, but weakness, numbness, trouble walking, or severe pain after major trauma needs urgent assessment.

\synonyms
Pain Whiplash (Whiplash)
Sprain Neck (Whiplash)
Strain Neck (Whiplash)


\medterm{Whiplash Injury} Alternate index label; see \textit{Whiplash}.

\medterm{Whipple Disease} A rare bacterial infection that can damage the small intestine, joints, brain, heart, or other organs. It may cause years of joint pain followed by weight loss, persistent diarrhea, abdominal pain, and poor nutrient absorption.

\medterm{Whipple Procedure} A major operation that removes the head of the pancreas and nearby parts of the small intestine, bile duct, and gallbladder, then reconnects the digestive tract. It is used for selected cancers and other serious diseases in this area.

\medterm{Whipple Triad} Three findings supporting that symptoms are caused by low blood sugar: typical symptoms occur, blood sugar is low at that time, and symptoms improve after blood sugar is raised.

\medterm{Whipple'S Disease} A rare bacterial infection caused by Tropheryma whipplei, damaging intestinal absorption and potentially affecting joints, brain, or heart.

\synonyms
Intestinal Lipodystrophy
Lipodystrophy, Intestinal

\medterm{Whipworm} A parasitic roundworm of the genus Trichuris, usually Trichuris trichiura in humans. Infection occurs by ingesting embryonated eggs and may cause abdominal pain, diarrhea, anemia, or rectal prolapse when severe.

\medterm{Whipworm (Trichuriasis)} Intestinal infection caused by the soil-transmitted helminth Trichuris trichiura
(whipworm), endemic in tropical and subtropical regions with poor sanitation. Humans
acquire infection by ingesting embryonated eggs from contaminated soil, food, or water.

\synonyms
Trichuriasis

\medterm{Whipworm Infection} Intestinal infection with Trichuris trichiura acquired by ingesting eggs from contaminated soil or food; heavy infection can cause abdominal symptoms, anemia, or rectal prolapse.

\medterm{Whirlpool Bath} Hydrotherapy device providing agitated water for therapeutic bathing; used in wound care, physical therapy, and rehabilitation.

\medterm{Whispered Pectoriloquy} An examination finding in which whispered words sound unusually clear through a stethoscope over the lung. It may suggest that air-filled lung has become solid, as in pneumonia, but it is not diagnostic alone.

\medterm{White Blood Cell} A cell of the immune system that helps protect the body from infections, foreign substances, and abnormal cells. White blood cells are made mainly in the bone marrow and include neutrophils, lymphocytes, monocytes, eosinophils, and basophils. Their number or function may change with infection, inflammation, immune disorders, or bone-marrow disease.

\medterm{White Blood Cell Count} A blood test measuring the number of white blood cells. A high or low result can occur with infection, inflammation, medicines, immune disorders, or bone-marrow disease and must be interpreted with the differential count and clinical findings.


\medterm{White Blood Cell Labeling} A scan method in which a person's white blood cells are removed, marked with a small amount of radioactive material, and returned to the body. The cells may collect at areas of infection or inflammation, which can then be seen on images.


\medterm{White Cell} Another name for a white blood cell, an immune cell made mainly in the bone marrow that helps defend the body against infection, foreign substances, and abnormal cells.

\medterm{White Cell Count} A blood test measuring the number of white blood cells and often the proportion of each type. Results help assess infection, inflammation, immune disorders, and bone-marrow problems.

\medterm{White Finger} Episodic pallor or whitening of the fingers caused by vasospasm, commonly associated with Raynaud phenomenon or hand-arm vibration syndrome.

\medterm{White Infarction} Tissue death caused by loss of blood flow in a solid organ such as the heart, kidney, or spleen. The damaged area looks pale because little blood can reach it through alternate vessels.

\synonyms
anemic infarction, occurs in solid organs with.

\medterm{White Leg (Phlegmasia Alba Dolens)} Painful, marked swelling and pallor of a limb caused by extensive deep venous thrombosis, usually involving major venous outflow. It requires urgent assessment.

\synonyms
Phlegmasia Alba Dolens

\medterm{White Matter} Tissue in the brain and spinal cord made mainly of nerve fibers covered by a protective substance called myelin. It carries signals between different brain areas and between the brain and spinal cord.

\medterm{White Matter Of The Brain} The brain tissue composed mainly of myelinated axons and supporting glial cells, connecting different brain regions and the brain with the spinal cord; disease or injury can impair cognition, sensation, coordination, or movement.

\medterm{White Of The Eye} The sclera, the tough, opaque outer coat of the eye surrounding the cornea and helping maintain the shape of the globe.

\medterm{White Spot Disease} Alternate index label; see \textit{Lichen sclerosus}.

\medterm{White Tongue} A white coating or patches on the tongue caused by debris, dry mouth, irritation, infection, inflammation, or other oral conditions.

\medterm{White-Coat Hypertension} Blood pressure that is high in a medical setting but normal outside it. Home or ambulatory readings help confirm the pattern, which may still predict higher future hypertension risk.

\synonyms
White-coat effect

\medterm{Whiteheads} Whiteheads are closed comedones formed when oil and dead skin block a hair follicle, producing small pale or skin-colored acne bumps.

\medterm{Whitlow} An infection or inflammation of the tissues around a finger or toenail, often involving the nail fold; it may be bacterial, fungal, or herpetic.

\medterm{Whitlow (Paronychia)} Acute or chronic infection and inflammation of the tissues around a fingernail or toenail, usually caused by bacteria, fungi, or repeated moisture and trauma.

\synonyms
Paronychia

\medterm{Whole Breast Radiation Therapy} Whole breast radiation therapy treats remaining breast tissue after breast-conserving surgery to reduce local recurrence risk, with planned dosing and monitoring for skin and tissue effects.

\medterm{Whole Exome Sequencing} A genetic test that reads most protein-coding sections of a person’s DNA. It can identify changes causing some inherited disorders, but it may miss changes outside coding regions, structural changes, or causes not yet understood; results require clinical interpretation.

\medterm{Whole Genome Sequencing} Laboratory analysis of nearly all of a person's nuclear DNA, including coding and noncoding regions, and sometimes mitochondrial DNA. It can identify many types of genetic variants but may produce uncertain or incidental findings and does not detect every cause of disease; results require appropriate clinical interpretation.

\medterm{Whole-Body Irradiation} Radiation treatment directed at the entire body, most often given before a stem-cell transplant to suppress the immune system and make room for donated cells. It can affect the bone marrow, fertility, and other organs.

\medterm{Whole-Body Scan} An imaging examination that surveys much or all of the body, using a technique such as PET/CT, CT, MRI, or a nuclear medicine scan. It may be used for selected cancer staging, infection localization, or evaluation of systemic disease; the coverage and usefulness depend on the test and clinical question.

\medterm{Whole-Bowel Irrigation} A hospital treatment that flushes the intestines with a special polyethylene-glycol drink. It is used only for selected poisonings or swallowed objects after clinicians assess the airway and risks.

\medterm{Whole-Brain Radiation Therapy} Radiation treatment directed at the entire brain, used in selected people with multiple brain metastases or cancer cells around the brain. It may also be used to reduce brain-spread risk in selected cancers; side effects depend on dose and treatment plan.

\medterm{Whole-Genome Sequencing} Sequencing of nearly the complete nuclear genome, including coding and noncoding
regions, together with mitochondrial DNA when covered by the assay. It can detect a
broad range of variants but generates substantial interpretation and incidental-finding
challenges.

\medterm{Whooping Cough} A highly contagious respiratory infection caused by \textit{Bordetella pertussis}. It causes prolonged bouts of rapid coughing that may be followed by a whooping sound or vomiting; infants may have pauses in breathing without the typical whoop. Also called pertussis.

\medterm{Whooping-Cough} Alternate index label; see \textit{Whooping Cough}.

\medterm{Widal Reaction} An older blood test for antibodies associated with typhoid or paratyphoid fever. It is not very reliable because results may reflect past infection or vaccination, so culture or molecular testing is usually preferred.


\medterm{Widely Spaced Teeth} Teeth separated by unusually large gaps, also called diastema when the gap is localized. Causes include normal development, missing teeth, tooth-size differences, habits, or jaw and gum conditions.

\medterm{Williams Syndrome} A genetic condition caused by loss of a small section of chromosome 7, including a gene important for elastic tissues. It may cause characteristic facial features, heart-vessel narrowing, high blood calcium, developmental differences, learning difficulties, and connective-tissue problems. Regular heart, calcium, hearing, and developmental checks are important.

\medterm{Wilm'S Tumour} Source spelling variant of \textit{Wilms tumor}, a malignant embryonal kidney tumor occurring mainly in young children.

\medterm{Wilms Tumor} A malignant embryonal tumor of the kidney in children. It may present as an abdominal mass, pain,
hematuria, or an incidental finding and may be associated with syndromes such as WAGR or
Beckwith-Wiedemann syndrome. Treatment and prognosis depend on stage and tumor biology.

\synonyms
Cancer, Wilms' Tumor

\medterm{Wilms' Tumor (Nephroblastoma)} An embryonal malignant tumor of the kidney that usually occurs in young children. It may be
associated with congenital syndromes, and treatment and prognosis depend on tumor stage,
histology, and molecular features.

\synonyms
Nephroblastoma

\medterm{Wilson Disease} An inherited disorder in which excess copper builds up in the liver, brain, eyes, and other organs. It can cause liver disease, tremor or other movement problems, mood or behavior changes, and colored rings around the cornea. Lifelong treatment removes or limits copper and prevents further buildup.

\medterm{Wilson'S Disease} An inherited condition in which copper builds up in the liver, brain, eyes, and other organs. It may cause liver disease, movement or mood changes, and colored rings around the cornea, but lifelong treatment can prevent further damage.

\medterm{Windowing} Adjusting the brightness and contrast of a CT image so particular tissues, such as bone, lung, or soft tissue, are easier to see.

\medterm{Windpipe} The tube carrying air from the voice box to the two main airways leading to the lungs. It is supported by firm rings that help keep it open.


\medterm{Winter Vomiting Disease} An acute gastroenteritis, usually caused by norovirus, characterized by sudden vomiting, diarrhea, and abdominal cramps.

\medterm{Wireless Motility Capsule} A swallowed capsule that records pressure, acidity, temperature, and movement as it travels through the digestive tract. The measurements estimate how long food takes to pass through the stomach and intestines and may help investigate unexplained constipation, diarrhea, nausea, or bloating.

\medterm{Wisdom Teeth} Wisdom teeth are the third molars that commonly emerge in late adolescence or adulthood. Impaction, infection, decay, cysts, gum disease, or crowding can cause pain and swelling; dental examination and imaging determine whether monitoring or removal is appropriate.

\medterm{Wisdom Teeth, Impacted} Alternate index label; see \textit{Impacted wisdom teeth}.

\medterm{Wisdom Tooth (Third Molar)} The last molar in each dental quadrant, usually erupting in late adolescence or early adulthood. Impaction may cause pain, infection, or damage to adjacent teeth.

\synonyms
Third Molar

\medterm{Wiskott-Aldrich Syndrome} A rare inherited immune disorder that mainly affects boys and causes eczema, repeated infections, and a low number of unusually small platelets. It can also cause bleeding and autoimmune problems and requires specialist care.

\medterm{Witch'S Milk} A temporary milky breast discharge in a newborn caused by maternal hormones crossing the placenta; it usually resolves without squeezing or treatment, which can otherwise cause infection.

\medterm{Witch'S Milk (Neonatal Galactorrhea)} A temporary milk-like breast secretion in a newborn caused by exposure to maternal hormones. It usually resolves without treatment; the breasts should not be squeezed.

\synonyms
Neonatal Galactorrhea

\medterm{Withdrawal} Physical or psychological symptoms that occur when a person who has developed dependence reduces or stops using a medicine or psychoactive substance. Symptoms vary by substance and may include anxiety, agitation, sweating, tremor, nausea, pain, insomnia, cravings, or seizures; some forms require urgent medical care.

\synonyms
Withdrawal syndrome

\medterm{Withdrawal Reflex} An automatic response in which a painful stimulus makes a limb pull away from danger. The opposite limb may straighten to help maintain balance.

\medterm{Withdrawal Symptoms} Physical or psychological symptoms that occur when a person stops or reduces a substance or medicine
after the body has adapted to it. The symptoms depend on the substance, duration of use,
and the person’s health.

\medterm{Withdrawal Syndrome} A predictable set of signs and symptoms that occur when a drug is discontinued or the
dose is reduced after chronic use, reflecting the physiological rebound from drug
adaptation.

\medterm{Wolf-Hirschhorn Syndrome} A rare genetic disorder caused by loss of part of chromosome 4. It can cause poor growth, a small head, distinctive facial features, developmental and intellectual disability, seizures, and heart or other organ abnormalities; lifelong specialist support is usually needed.

\synonyms
4p- Syndrome; 4p Deletion Syndrome

\medterm{Wolff Parkinson White Syndrome (Wolff-Parkinson-White Syndrome)} Alternate index label for Wolff-Parkinson-White Syndrome.

\medterm{Wolff-Parkinson-White Accessory Pathway-Related Tachycardia} Alternate index label; see \textit{Wolff-Parkinson-White syndrome}.

\medterm{Wolff-Parkinson-White Syndrome} A heart condition present from birth in which an extra electrical pathway can make the heart beat very fast. It may cause sudden palpitations, dizziness, fainting, or rarely cardiac arrest and can be treated by medicines or a procedure.

\synonyms
Wolff Parkinson White Syndrome (Wolff-Parkinson-White Syndrome)

\medterm{Wolffian Ducts} Temporary tubes in the developing embryo that form parts of the male reproductive system, including the sperm-transport ducts. They usually shrink during typical female development, although small remnants may remain.

\medterm{Wolfram Syndrome} A rare inherited disorder that commonly causes diabetes requiring insulin and progressive loss of vision from optic-nerve damage. It may also cause excessive thirst and urination, hearing loss, bladder problems, balance difficulty, or other nerve-system problems. Treatment monitors and supports each affected function.

\medterm{Wolinella} Genus of anaerobic, Gram-negative bacteria found in the oral cavity and gastrointestinal tract; associated with periodontal disease.


\medterm{Wolman Disease} A severe inherited disorder in which a missing enzyme causes fats to build up in organs. It usually begins in infancy with vomiting, diarrhea, poor growth, an enlarged liver and spleen, and adrenal-gland problems. Without treatment it can be fatal early; enzyme replacement and specialist care are required.

\medterm{Wolman'S Disease} A severe lysosomal acid-lipase deficiency presenting in infancy with hepatosplenomegaly,
adrenal calcification, malabsorption, vomiting, diarrhea, and failure to thrive. It is a
severe form of lysosomal acid-lipase deficiency and can be fatal early without
treatment.

\medterm{Womb} Common name for the uterus; the muscular organ where a fetus develops during pregnancy.
\synonyms
Uterus

\medterm{Womb Growths (Uterine Fibroids)} Alternate index label for Uterine Fibroids.

\medterm{Women'S Health} Medical specialty focusing on conditions unique to or more prevalent in females, including reproductive health, pregnancy, menopause, and gender-specific disease patterns.

\synonyms
Womens Health (Women's Health)


\medterm{Womens Health (Women'S Health)} Alternate index label for Women's Health.

\medterm{Wood Lamp Examination} Examination of the skin in a dark room using special ultraviolet light. The light can make some color changes and infections easier to see, but the finding usually needs confirmation with the examination and other tests.


\medterm{Wood Spirit} An old name for methanol, a toxic alcohol whose ingestion can cause metabolic acidosis, blindness, and death.

\medterm{Wood'S Light} A special ultraviolet light used in a darkened room to make some skin colors, scales, and infections easier to see. It can support an examination but usually cannot identify the cause by itself.

\medterm{Woods Corkscrew Maneuver} An emergency procedure used when a baby’s shoulders become stuck during birth. A trained clinician rotates the baby’s shoulders to help release them and complete delivery.

\medterm{Woods Screw Maneuver} A rotational maneuver used during shoulder dystocia to rotate the fetal shoulders and assist delivery.
\medterm{Woody Tongue} A rare long-lasting infection, usually caused by \textit{Actinomyces}, that makes the tongue and nearby tissues hard and swollen. It can interfere with eating, speaking, or breathing.

\medterm{Wool-Sorter'S Disease (Anthrax)} A form of inhaled anthrax acquired by breathing spores of \textit{Bacillus anthracis}, historically associated with handling contaminated wool or animal products. It can cause severe breathing illness and requires urgent treatment.

\synonyms
Anthrax


\medterm{Word Blindness} An older term for alexia, an acquired inability to read or recognize written words despite adequate vision and previously developed language ability, usually caused by brain dysfunction or injury.

\medterm{Word Blindness (Alexia)} An acquired neurological disorder in which a person loses or markedly impairs the ability to read despite adequate vision and previously acquired language skills.

\synonyms
Alexia

\medterm{Word Deafness (Auditory Verbal Agnosia)} An acquired inability to understand spoken words even though ordinary hearing is intact. It usually results from damage to brain areas that connect hearing with language, such as after a stroke or head injury.

\synonyms
Auditory Verbal Agnosia

\medterm{Word Recognition Scoring} A hearing test measuring how accurately a person repeats or identifies spoken words at a comfortable volume. It helps assess how well the hearing system and brain understand speech.

\medterm{Word Salad} Speech made of words or phrases that have little or no logical connection, making the person difficult to understand. It can occur with severe psychosis, mania, brain disease, or intoxication and needs assessment in context.

\medterm{Working Memory} The ability to hold a small amount of information briefly and use it while completing a task, such as remembering instructions while following them. Attention, stress, sleep, illness, and brain conditions can affect it.

\medterm{Worm Infestation (Helminthiasis)} Infection by parasitic worms, including roundworms, tapeworms, and flukes. Symptoms depend on the species and body site and may include abdominal problems, anemia, cough, itching, or no symptoms.

\synonyms
Helminthiasis

\medterm{Worms Pinworms (Pinworm Infection)} Alternate index label for Pinworm Infection.

\medterm{Wound} Damage to body tissue caused by trauma, surgery, pressure, heat, chemicals, or another cause. A wound may be open or closed and can vary in depth, contamination, bleeding, and injury to nearby nerves, vessels, or organs.

\medterm{Wound Care} Assessment and treatment of a wound through cleansing, bleeding control, removal of dead or contaminated tissue when indicated, appropriate closure or dressing, pain control, and follow-up. Care also includes tetanus prevention, pressure relief, and treatment of infection or underlying conditions; antibiotics are reserved for selected contaminated wounds, bites, or infections.

\medterm{Wound Complication} A problem that develops during healing, such as infection, separation of wound edges, bleeding, fluid
collection, tissue death, or excessive scarring. Increasing pain, redness, swelling, drainage,
fever, or wound opening should be assessed promptly.

\medterm{Wound Dehiscence} Partial or complete reopening of a surgical wound after it has been closed. It may follow infection, poor blood supply, diabetes, increased pressure, or strain on the wound. New separation, drainage, fever, severe pain, or visible deeper tissue needs prompt medical assessment.

\medterm{Wound Evisceration} Separation of a surgical wound with protrusion of internal organs or tissue through the opening. It is a surgical emergency requiring protection of the exposed tissue and urgent medical evaluation.

\medterm{Wound Healing} The body’s process of stopping bleeding, controlling inflammation, replacing damaged tissue, covering the surface, and strengthening the repair. Blood supply, infection, nutrition, diabetes, medicines, smoking, wound depth, and movement can affect healing.

\medterm{Wound Healing Phases} Wound healing has overlapping stages: stopping bleeding, inflammation, building new tissue and skin, and strengthening the repair. The timing and result depend on the wound, blood supply, infection, nutrition, and general health.

\medterm{Wound Infection} Infection of a wound by microorganisms, causing findings such as increasing pain, redness, warmth,
swelling, pus, fever, or delayed healing.

\medterm{Wound Management} Assessment and treatment of acute or chronic wounds through cleansing, debridement when indicated, moisture and exudate control, infection management, pressure relief, nutrition, and correction of underlying disease.
\medterm{Wound VAC} A negative-pressure wound treatment in which a sealed dressing is connected to a pump that gently removes fluid and reduces swelling. It can support new tissue growth in selected complex wounds, but requires proper application and monitoring for bleeding, infection, or skin damage.

\medterm{Wounds} Injuries in which tissue is damaged by trauma, surgery, pressure, heat, chemicals, or other causes; wounds may be open or closed and may require cleansing, closure, or debridement.

\medterm{Wright Stain} A laboratory dye applied to a blood sample so different blood cells and abnormal cell changes can be seen and identified under a microscope.



\medterm{Wrisberg Ganglion} A small cluster of nerve cells connected with the facial nerve. It helps relay taste signals from part of the tongue and sensory signals from nearby areas of the ear.

\medterm{Wrist} The joint connecting the hand to the forearm. It contains eight small wrist bones, ligaments, tendons, nerves, and blood vessels and allows the hand to bend, turn, and move from side to side.

\medterm{Wrist Arthroscopy} A procedure using a small camera and instruments inserted through tiny wrist openings to examine or treat damaged cartilage, ligaments, joint lining, or other wrist problems.

\medterm{Wrist Drop} Inability to lift the wrist and fingers, usually caused by injury or compression of the radial nerve. It may cause weakness, numbness, or difficulty using the hand and needs assessment for the cause.
\medterm{Wrist Fracture} A break in one or more bones near the wrist, most commonly the distal radius or ulna, usually caused by a fall, direct trauma, or weakened bone.

\medterm{Wrist Injury} Damage to the bones, ligaments, tendons, nerves, or vessels of the wrist from a fall, twist, impact, or overuse; fracture, instability, severe swelling, or neurovascular symptoms require assessment.

\medterm{Wrist Joint} The joint where the forearm bones meet the small bones of the hand. It allows bending, straightening, side-to-side movement, and limited rotation while ligaments and cartilage provide stability.

\medterm{Wrist Pain} Pain in or around the wrist, commonly caused by sprain, fracture, tendinitis, arthritis,
or nerve compression such as carpal tunnel syndrome. Persistent pain, deformity,
numbness, or weakness requires assessment.

\synonyms
Pain, wrist

\synonyms
Pain, Wrist

\medterm{Wrist Sprain} Stretching or tearing of one or more ligaments supporting the wrist, causing pain, swelling, bruising, and reduced movement after injury.


\medterm{Writer'S Cramp} A movement disorder in which involuntary muscle tightening causes abnormal hand or arm positions while writing. It may also affect other closely related fine hand movements.

\medterm{Written Documentation} Formal recording of patient information, clinical decisions, and treatment plans; essential for continuity of care and medicolegal purposes.

\medterm{Wry-Neck} Torticollis, an abnormal tilt or rotation of the head caused by involuntary contraction, shortening, or dysfunction of neck muscles.

\medterm{Wryneck (Torticollis)} Abnormal tilting or rotation of the head caused by involuntary contraction, shortening, or dysfunction of neck muscles. It may be congenital, acquired, painful, or neurological.

\medterm{Wt} A medical abbreviation commonly meaning weight. Its meaning should be confirmed from the surrounding documentation because abbreviations can be ambiguous.

\medterm{Wuchereria} A genus of filarial nematodes that includes Wuchereria bancrofti, a major cause of
lymphatic filariasis. Infection is transmitted by mosquitoes and may cause lymphangitis,
chronic lymphedema, hydrocele, and elephantiasis.

\medterm{Wuchereria Bancrofti} A parasitic roundworm spread by mosquitoes and the main cause of lymphatic filariasis. Long-term infection can block lymph drainage and cause swelling of the limbs or genitals, sometimes becoming severe and permanent.

\medterm{Wuhan Coronavirus (2019-NCoV)} Alternate index label for 2019-nCoV.

\medterm{Weight Loss Surgery} Weight-loss surgery, or bariatric surgery, changes the stomach or digestive tract to help treat severe obesity and related conditions. It requires long-term nutritional, medical, and lifestyle follow-up.

\medterm{Wheelchair} A wheelchair is a mobility device with a seat and wheels that helps a person move when walking is difficult or unsafe. Proper fitting and training improve comfort, safety, and independence.

\medterm{Wheelchairs} Wheelchairs are mobility devices used by people who cannot walk safely or efficiently. They may be manual or powered and should be selected and adjusted to the person’s posture, strength, environment, and daily needs.
